\documentclass[11pt, reqno]{amsart}

\input{~/latex-common/macros.tex}
\usepackage[backend=bibtex,style=science]{biblatex}
% \usepackage{tikzfill}
% \usetikzlibrary{fill}
% \bibliography{main.bib}
\pgfplotsset{compat=1.18}

\pagestyle{fancy}                         % fancy (allow headers, footers)
\fancyhf{}                                % clear all header/footer settings.
\cfoot{\thepage}                          % set page-numbers in footer.
% \lhead{\textit{\textbf{ Amittai, S}}}   % set name in header, left.
% \rhead{\textsc{Math 71: Algebra}}       % set class name in header, right.
\renewcommand{\headrulewidth}{0pt}
\renewcommand{\footrulewidth}{0pt}


\renewcommand{\theenumi}{\alph{enumi}}

\begin{document}

\newdate{due-date}{27}{04}{2023}

\title{CS-89.31: Deep Learning Generalization and Robustness\\ Amittai Siavava \\ \displaydate{due-date}}
\author{Amittai Siavava}
% \date{\today}


\setlength{\headheight}{13.0pt}
\setlength{\footskip}{15.0pt}

\maketitle

\bigskip

\def \cram { \textsc{cram} }
\def \dom { \textsc{domineering} }
\def \sub { \textsc{subtraction} }
\def \weighted { \textsc{weighted odds and evens}}
\def \nim { \textsc{nim} }
\def \P { \mathbf{P}}
\def \N { \mathbf{N}}

\newpage
\begin{problem}[1]
  Let us denote with:
  \begin{itemize}
    \item $N \times N\times V$ the size of the input volume;
    \item $n \times n \times V$ the size of the filters;
    \item $Z$ the number of zeros padded at the top/bottom/left/right of the image;
    \item $S$ the stride;
    \item $V'$ the number of filters.
  \end{itemize}
  Prove the output volume will have size $M \times M \times V'$,
  where $M = \frac{N - n+2Z}{S} + 1$.

  \step
  \emph{
    [Hint] Drawing the convolution operation with $Z$ and $S$ on it
    can help you get some intuition of proving the formula.
  }

  \step
  \begin{enumarabic}
    \item Since we pad the input volume with $Z$ zeros on each side,
      the resulting volume (before convolution) is $N + 2Z$.
      The resulting volume retains the same depth $V$.
    \item Next, we apply filters of size $n \times n \times V$.
      Since filters cannot be partially applied(i.e. part of the filter being in the
      image and part of it being outside), we can only apply the filter
      from position $0$ to position $(N + 2Z) - n$ (when the right edge of the filter
      will be at position $N + 2Z$).
    \item A stride $S$ subdivides this interval into $\frac{(N + 2Z) - n}{S}$ intervals.
      However, since both the starting position and the ending position are
      applied to, we get an extra position for the filter, making
      the number of intervals $\frac{(N + 2Z) - n}{S} + 1$.
    \item Therefore, each applied filter will produce a volume of size
      $M \times M$, where $M = \frac{(N + 2Z) - n}{S} + 1$.
      Since we have $V'$ filters, the resulting volume will be of size
      $M \times M \times V'$.

      \step
      \emph{See next page for drawing.}
  \end{enumarabic}

  \newpage
  For simplicity, I fixed the volume to $9 \times 9 \times 1$, and
  also fixed the filter to $3 \times 3 \times 1$.
  In direct contrast to the formula above, $(N + 2Z) = 9$ and $n = 3$.

      \noindent
      \begin{figure}[h]
        \begin{minipage}{.3\textwidth}
          \centering
          \begin{tikzpicture}[scale=0.3]
            % draw a 9X9 grid
            \draw[step=1cm,gray,very thin] (0, 0) grid (9, 9);
          \end{tikzpicture}
          \caption{\\Input-volume $(9~\times~9)$}
        \end{minipage}%
        \begin{minipage}{.5\textwidth}
          \centering
          \begin{tikzpicture}[scale=0.3]
            \fill[electriclime] (0,0) rectangle (3, 3);
            \draw[step=1cm,black,very thin] (0,0) grid (3,3);
          \end{tikzpicture}
          \caption{Filter}
        \end{minipage}

        \step
        \begin{minipage}{.3\textwidth}
          \centering
          \begin{tikzpicture}[scale=0.3]
            % draw filter at position 0
            \fill[electriclime] (0, 9) rectangle (3, 6);
            \fill[glaucous] (1, 8) rectangle (4, 5);
            \fill[electriclime] (2, 7) rectangle (5, 4);
            \fill[glaucous] (3, 6) rectangle (6, 3);
            \fill[electriclime] (4, 5) rectangle (7, 2);
            \fill[glaucous] (5, 4) rectangle (8, 1);
            \fill[electriclime] (6, 3) rectangle (9, 0);
            % draw a 9X9 grid
            \draw[step=1cm,black,very thin] (0,0) grid (9,9);
          \end{tikzpicture}
          \caption{\\Stride~=~$1$\\ $M = 7 = \frac{9 - 3}{1} + 1$}
        \end{minipage}%
        \begin{minipage}{.3\textwidth}
          \centering
          \begin{tikzpicture}[scale=0.3]
            % draw filter at position 0
            \fill[electriclime] (0,9) rectangle (3, 6);
            \fill[glaucous] (2, 7) rectangle (5, 4);
            \fill[electriclime] (4, 5) rectangle (7, 2);
            \fill[glaucous] (6, 3) rectangle (9, 0);
            % draw a 9X9 grid
            \draw[step=1cm,black,very thin] (0,0) grid (9,9);
          \end{tikzpicture}
          \caption{\\Stride~=~$2$\\ $M = 4 = \frac{9 - 3}{2} + 1$}
        \end{minipage}%
        \begin{minipage}{.3\textwidth}
          \centering
          \begin{tikzpicture}[scale=0.3]
            % draw filter at position 0
            \fill[electriclime] (0, 9) rectangle (3, 6);
            \fill[glaucous] (3, 6) rectangle (6, 3);
            \fill[electriclime] (6, 3) rectangle (9, 0);
            % draw a 9X9 grid
            \draw[step=1cm,black,very thin] (0,0) grid (9,9);
          \end{tikzpicture}
          \caption{\\Stride~=~$3$\\ $M = 3 = \frac{9 - 3}{3} + 1$}
        \end{minipage}%
      \end{figure}
  
\end{problem}
\end{document}
